\begin{helper}
Assume
\(\noderef{ParameterHierarchyEventualComparisons}
(\varepsilon,\varepsilon_1,\varepsilon_2,n_0)\), \(n>n_0\),
\[
m=\noderef{TwoBiteNaturalReducedVertexCount}(n),
\qquad
p=\noderef{TwoBiteEdgeProbability}\!\left(\frac12,n\right),
\]
and \(\noderef{FiberAndDegreeEvent}(\omega,\varepsilon_2)\).  Then every red
base vertex \(r\) and every blue base vertex \(b\) satisfy
\[
\left|
  X_{\operatorname{inl}(r)}(I).\mathrm{image}
    (\noderef{TwoBiteRedProjection}(\omega))
\right|
\le 2pm,
\]
and
\[
\left|
  X_{\operatorname{inr}(b)}(I).\mathrm{image}
    (\noderef{TwoBiteBlueProjection}(\omega))
\right|
\le 2pm.
\]
\end{helper}
\begin{proof}
We prove the red assertion first.  Fix a red base vertex \(r\).  If
\[
u\in X_{\operatorname{inl}(r)}(I).\mathrm{image}
  (\noderef{TwoBiteRedProjection}(\omega)),
\]
then by the definition of finite-set image there is some
\(v\in X_{\operatorname{inl}(r)}(I)\) with
\(u=\noderef{TwoBiteRedProjection}(\omega,v)\).  Unfolding
\(\noderef{TwoBiteX}\) and the red branch of
\(\noderef{TwoBiteLiftedNeighborhood}\), membership of \(v\) in
\(X_{\operatorname{inl}(r)}(I)\) gives
\[
(\noderef{TwoBiteRedGraph}(\omega)).Adj
  \ r\ (\noderef{TwoBiteRedProjection}(\omega,v)).
\]
Substituting \(u\), this says that \(u\) lies in the filtered-universe red
neighbor set
\[
\{w:(\noderef{TwoBiteRedGraph}(\omega)).Adj\ r\ w\}.
\]
Therefore the red projected image is contained in this neighbor set, and
finite-cardinality monotonicity gives an upper bound by its cardinality.  By
the definition of \(\noderef{GraphDegreeCount}\), that cardinality is exactly
\[
\noderef{GraphDegreeCount}(\noderef{TwoBiteRedGraph}(\omega),r).
\]

The red degree clause of \(\noderef{FiberAndDegreeEvent}\) gives
\[
\left|
  \noderef{GraphDegreeCount}(\noderef{TwoBiteRedGraph}(\omega),r)-pm
\right|\le \varepsilon_2pm.
\]
Hence the upper side of the absolute-value inequality gives
\[
\noderef{GraphDegreeCount}(\noderef{TwoBiteRedGraph}(\omega),r)
\le (1+\varepsilon_2)pm.
\]
The eventual-comparisons package, with the displayed identifications of
\(m\) and \(p\), gives \(1\le2pm\) and \(\varepsilon_2\le1\).  Thus
\(pm\ge0\), and multiplying \(\varepsilon_2\le1\) by \(pm\) gives
\[
(1+\varepsilon_2)pm\le2pm.
\]
Combining the image containment, degree estimate, and this last inequality
proves the red bound.

The blue proof is identical with colors exchanged.  For
\(u\in X_{\operatorname{inr}(b)}(I).\mathrm{image}
(\noderef{TwoBiteBlueProjection}(\omega))\), choose
\(v\in X_{\operatorname{inr}(b)}(I)\) mapping to \(u\).  Unfolding
\(\noderef{TwoBiteX}\) and the blue branch of
\(\noderef{TwoBiteLiftedNeighborhood}\) gives adjacency from \(b\) to \(u\) in
\(\noderef{TwoBiteBlueGraph}(\omega)\).  Hence the image is contained in the
filtered-universe blue neighbor set, whose cardinality is the corresponding
\(\noderef{GraphDegreeCount}\).  The blue degree clause of
\(\noderef{FiberAndDegreeEvent}\), followed by the same
\(\varepsilon_2\le1\) and \(pm\ge0\) calculation, gives the bound \(2pm\).
\end{proof}
